\documentclass[11pt]{article}
\usepackage[margin=1in]{geometry}
\usepackage{enumitem}
% =============================================================================
% Preambles/header.tex — shared LaTeX/Beamer preamble for this project
%
% Use from a Beamer lecture by prepending:
%     \input{header}
% and compiling with TEXINPUTS=../Preambles:$TEXINPUTS (the /compile-latex
% skill does this automatically).
%
% PALETTE CONTRACT. Color names defined here MUST match the names in
% Quarto/theme-template.scss so Beamer and Quarto renderings use the same
% palette. The scripts/check-palette-sync.sh script enforces this.
%
% To customize: edit the HEX values below AND the matching $variable in
% Quarto/theme-template.scss, then run scripts/check-palette-sync.sh.
% =============================================================================

% -----------------------------------------------------------------------------
% Engine detection + encoding
%
% Our /compile-latex skill uses XeLaTeX, where inputenc is unnecessary and
% can produce encoding warnings. Only load inputenc under pdfTeX.
% -----------------------------------------------------------------------------
\usepackage{iftex}
\ifPDFTeX
  \usepackage[utf8]{inputenc}
\fi

\usepackage{amsmath, amssymb, amsthm}
\usepackage{booktabs}
\usepackage{graphicx}

% -----------------------------------------------------------------------------
% PALETTE — mirrors Quarto/theme-template.scss variable names.
% Load xcolor and define colors BEFORE anything that references them
% (notably hyperref, hypersetup, and the Beamer color assignments).
% -----------------------------------------------------------------------------
\usepackage{xcolor}

\definecolor{primary-blue}{HTML}{012169}      % headings, accents
\definecolor{primary-gold}{HTML}{B9975B}      % emphasis, borders
\definecolor{highlight-yellow}{HTML}{F2A900}  % markers, alerts
\definecolor{light-bg}{HTML}{E8EDF5}          % title / section backgrounds
\definecolor{jet}{HTML}{1A1A1A}               % body text

% Semantic colors (used in diagrams and annotations)
\definecolor{positive}{HTML}{15803D}          % good / observed / identified
\definecolor{negative}{HTML}{B91C1C}          % bad / problematic / bias
\definecolor{neutral}{HTML}{525252}           % reference / context

% Highlight accents (match .hi-* classes in the SCSS)
\definecolor{hi-slate}{HTML}{314F4F}
\definecolor{hi-green}{HTML}{2E8B57}
\definecolor{hi-red}{HTML}{C41E3A}

% Semantic aliases so skill templates and snippets can write
% \textcolor{accent}{...} without hard-coding "primary-blue".
\colorlet{accent}{primary-blue}
\colorlet{accent2}{primary-gold}

% -----------------------------------------------------------------------------
% Hyperref — loaded AFTER xcolor + palette so link colors resolve.
% -----------------------------------------------------------------------------
\usepackage{hyperref}
\hypersetup{colorlinks=true, linkcolor=primary-blue, urlcolor=primary-blue,
            citecolor=primary-blue, pdfborder={0 0 0}}

% -----------------------------------------------------------------------------
% Course code — set once per deck (after \input{header}, before \begin{document})
% via \coursecode{CS401}. Shown in the footer on every slide so a deck's course
% is identifiable without opening the title page. Empty by default.
% -----------------------------------------------------------------------------
\makeatletter
\newcommand{\coursecode}[1]{\gdef\@coursecodetext{#1}}
\gdef\@coursecodetext{}
\makeatother

% -----------------------------------------------------------------------------
% Beamer theme assignments (only apply under Beamer).
%
% We detect Beamer via \@ifclassloaded{beamer}, which is the documented,
% non-internal way. This must be wrapped in \makeatletter/\makeatother
% because \@ifclassloaded contains an @-catcode character.
% -----------------------------------------------------------------------------
\makeatletter
\@ifclassloaded{beamer}{%
  \usefonttheme{professionalfonts}
  \setbeamercolor{structure}{fg=primary-blue}
  \setbeamercolor{frametitle}{fg=primary-blue}
  \setbeamercolor{title}{fg=primary-blue}
  \setbeamercolor{subtitle}{fg=primary-gold}
  \setbeamercolor{author}{fg=primary-blue}
  \setbeamercolor{institute}{fg=neutral}
  \setbeamercolor{date}{fg=primary-gold}
  \setbeamercolor{itemize item}{fg=primary-blue}
  \setbeamercolor{itemize subitem}{fg=primary-gold}
  \setbeamercolor{alerted text}{fg=primary-gold}
  \setbeamercolor{block title}{fg=white, bg=primary-blue}
  \setbeamercolor{block body}{bg=light-bg}
  % Semantic block variants: block=definition/concept (blue), exampleblock=
  % worked example (green/positive), alertblock=key takeaway or caution (gold).
  % Keep to <=2 colored boxes per slide across ALL three variants (INV-7).
  \setbeamercolor{block title example}{fg=white, bg=positive}
  \setbeamercolor{block body example}{bg=light-bg}
  \setbeamercolor{block title alerted}{fg=white, bg=primary-gold}
  \setbeamercolor{block body alerted}{bg=light-bg}

  % Minimal footer: course code (left) + page X / N (right), no navigation clutter
  \setbeamertemplate{navigation symbols}{}
  \setbeamertemplate{footline}{%
    \color{neutral}\footnotesize\hspace{0.7em}\@coursecodetext\hfill
    \insertframenumber/\inserttotalframenumber\hspace{0.7em}\vspace{0.3em}}
}{}
\makeatother

% -----------------------------------------------------------------------------
% TikZ setup — libraries the snippet gallery uses
% -----------------------------------------------------------------------------
\usepackage{tikz}
\usetikzlibrary{arrows.meta, positioning, calc, decorations.pathreplacing,
                fit, shapes.geometric, backgrounds}

% Shared TikZ styles — snippets and hand-written diagrams can pick these up
% via `\begin{tikzpicture}[every node/.style={...}]` or by naming specific
% styles (e.g., `\node[dag-node] (X) at (0,0) {X};`).
\tikzset{
  % Node styles for causal DAGs and similar
  dag-node/.style = {
    draw=primary-blue, fill=light-bg, rounded corners=2pt,
    minimum width=1.2cm, minimum height=0.9cm, align=center, font=\small
  },
  decision-node/.style = {
    draw=primary-gold, fill=white, diamond, aspect=1.8,
    minimum width=1.8cm, minimum height=1.2cm, align=center, font=\small,
    inner sep=1pt
  },
  % Edge styles — observed vs counterfactual
  observed-edge/.style = {-{Latex[length=2.5mm]}, thick, draw=jet},
  counterfactual-edge/.style = {-{Latex[length=2.5mm]}, thick, draw=neutral, dashed},
  confound-edge/.style = {-{Latex[length=2.5mm]}, thick, draw=negative, dashed},
  % Data-point styles for plots
  observed-dot/.style = {fill=primary-blue, draw=primary-blue, circle, inner sep=1.2pt},
  counterfactual-dot/.style = {fill=white, draw=neutral, circle, inner sep=1.2pt}
}

% -----------------------------------------------------------------------------
% Convenience macros
% -----------------------------------------------------------------------------
\newcommand{\muted}[1]{\textcolor{neutral}{#1}}
\newcommand{\key}[1]{\textcolor{primary-gold}{\textbf{#1}}}
\newcommand{\good}[1]{\textcolor{positive}{#1}}
\newcommand{\bad}[1]{\textcolor{negative}{#1}}

% Standout frame macro: full-bleed dark background for section transitions.
% Beamer-only (uses \begin{frame}/\setbeamercolor/\usebeamerfont) -- do not
% call from an article-class file (e.g. Notes/<CODE>/*-notes.tex). \muted,
% \key, \good, \bad, and \coursecode above are plain xcolor/gdef and are
% safe to use from either class; verified when Notes/article-class support
% was added.
% Expand: \transitionslide{Your transition title}
\newcommand{\transitionslide}[1]{%
  \begingroup
  \setbeamercolor{background canvas}{bg=primary-blue}
  \begin{frame}[plain]
    \centering
    \vfill
    {\usebeamerfont{title}\color{white}\Large #1\par}
    \vfill
  \end{frame}
  \endgroup
}

% Section divider frame (Beamer-only), an alternative to \transitionslide with a
% darker two-line style: near-black (jet) background, a small white label line
% above a larger gold title, and a thin gold rule along the bottom edge.
% Expand: \sectiondivider{Section 2}{Pointers}
\newcommand{\sectiondivider}[2]{%
  \begingroup
  \setbeamercolor{background canvas}{bg=jet}
  \begin{frame}[plain]
    \vspace*{3.0cm}
    {\color{white}\ttfamily\large #1\par}
    \vspace{0.5cm}
    {\color{primary-gold}\LARGE #2\par}
    \begin{tikzpicture}[remember picture, overlay]
      \draw[primary-gold, line width=2.5pt]
        ($(current page.south west)+(0,0.1)$) -- ($(current page.south east)+(0,0.1)$);
    \end{tikzpicture}
  \end{frame}
  \endgroup
}

% -----------------------------------------------------------------------------
% Channel watermark — "Concepts That Click" (YouTube).
%
% Article-class files (CompetitiveExam Q/A sets, Notes, Assignments) get a
% light diagonal draft watermark on every page plus a centered footer naming
% the channel. Beamer decks get a subtle TikZ background watermark instead
% (the footline already carries the course code). Centralized here so every
% MCQ PDF is branded without per-file edits.
% -----------------------------------------------------------------------------
\makeatletter
\@ifclassloaded{article}{%
  \usepackage{draftwatermark}
  \SetWatermarkText{Concepts That Click}
  \SetWatermarkScale{1}
  \SetWatermarkFontSize{2cm}
  \SetWatermarkLightness{0.86}
  \SetWatermarkAngle{45}
  \usepackage{fancyhdr}
  \pagestyle{fancy}
  \fancyhf{}
  \fancyfoot[C]{\footnotesize\textcolor{neutral}{Concepts That Click~|~YouTube: \href{https://www.youtube.com/@concepts.thatclick}{@concepts.thatclick}}}
  \renewcommand{\headrulewidth}{0pt}
  \renewcommand{\footrulewidth}{0pt}
  \fancypagestyle{plain}{%
    \fancyhf{}%
    \fancyfoot[C]{\footnotesize\textcolor{neutral}{Concepts That Click~|~YouTube: \href{https://www.youtube.com/@concepts.thatclick}{@concepts.thatclick}}}%
    \renewcommand{\headrulewidth}{0pt}%
    \renewcommand{\footrulewidth}{0pt}%
  }
}{}
\@ifclassloaded{beamer}{%
  \setbeamertemplate{background}{%
    \begin{tikzpicture}[remember picture, overlay]
      \node[rotate=45, opacity=0.055, align=center]
        at (current page.center)
        {\Huge\textcolor{neutral}{Concepts That Click}};
    \end{tikzpicture}%
  }
}{}
\makeatother 
\coursecode{JKSSB}
\title{Lesson 47 Practice: Humanware and IT Governance}
\author{JKSSB Computer Awareness}
\date{\today}
\begin{document}
\maketitle
\noindent\textbf{Provenance:} All 20 questions are original JKSSB-pattern
questions (\textit{[Original]}); none reproduces an official paper item.

\begin{enumerate}[label=\textbf{Q\arabic*.}, leftmargin=*]
\item \textit{[Original]} What does the term \textbf{humanware} refer to in a
  computer system?
  \begin{enumerate}[label=\Alph*.]
    \item The physical parts of the computer
    \item The programs and instructions
    \item The people involved with the computer system
    \item The firmware stored in the ROM
  \end{enumerate}
\item \textit{[Original]} Which of the following is \textbf{NOT} a category of
  humanware?
  \begin{enumerate}[label=\Alph*.]
    \item Users
    \item Operators
    \item System administrators
    \item Firmware
  \end{enumerate}
\item \textit{[Original]} Expand ICT.
  \begin{enumerate}[label=\Alph*.]
    \item Information and Computer Technology
    \item Information and Communication Technology
    \item Internet and Communication Technology
    \item Integrated Computer Technology
  \end{enumerate}
\item \textit{[Original]} Which component of a computer system consists of the
  people who design, operate and use it?
  \begin{enumerate}[label=\Alph*.]
    \item Hardware
    \item Software
    \item Humanware
    \item Firmware
  \end{enumerate}
\item \textit{[Original]} Which statement correctly distinguishes
  \textbf{IT governance} from \textbf{e-governance}?
  \begin{enumerate}[label=\Alph*.]
    \item IT governance is governance \emph{of} IT; e-governance is IT
      \emph{in} governance.
    \item IT governance is IT \emph{in} governance; e-governance is
      governance \emph{of} IT.
    \item They mean exactly the same thing.
    \item IT governance is a type of computer hardware.
  \end{enumerate}
\item \textit{[Original]} E-governance refers to:
  \begin{enumerate}[label=\Alph*.]
    \item the use of ICT by government to deliver services and interact with
      citizens.
    \item the management of a company's internal IT resources.
    \item the repair of government computers.
    \item the sale of software to the government.
  \end{enumerate}
\item \textit{[Original]} In e-governance, G2C stands for:
  \begin{enumerate}[label=\Alph*.]
    \item Government to Company
    \item Government to Citizen
    \item Government to Community
    \item Government to Court
  \end{enumerate}
\item \textit{[Original]} Which pairing of an e-governance interaction type
  with its meaning is correctly matched?
  \begin{enumerate}[label=\Alph*.]
    \item G2B --- Government to Employee
    \item G2G --- Government to Citizen
    \item G2E --- Government to Employee
    \item G2C --- Government to Business
  \end{enumerate}
\item \textit{[Original]} Which of the following is a benefit of using IT in
  governance?
  \begin{enumerate}[label=\Alph*.]
    \item Increased transparency and accountability
    \item Increased use of paperwork
    \item Slower processing of applications
    \item Greater need for citizens to visit offices in person
  \end{enumerate}
\item \textit{[Original]} Evaluate the statements:
  \begin{enumerate}[label=\Roman*.]
    \item E-governance uses ICT to deliver government services to citizens.
    \item IT governance means using IT to govern the state.
    \item G2G means Government to Government.
  \end{enumerate}
  \begin{enumerate}[label=\Alph*.]
    \item I only
    \item I and II only
    \item I and III only
    \item I, II and III
  \end{enumerate}
\item \textit{[Original]} An \textbf{electronic record} is:
  \begin{enumerate}[label=\Alph*.]
    \item a printed paper file kept in a government office.
    \item information created, stored and retrieved in electronic form.
    \item a hardware device for storing paper.
    \item a type of computer virus.
  \end{enumerate}
\item \textit{[Original]} Which of the following is an example of an
  electronic record?
  \begin{enumerate}[label=\Alph*.]
    \item A scanned land document stored on a computer
    \item A typed letter kept only on paper
    \item A printed certificate
    \item A handwritten office register
  \end{enumerate}
\item \textit{[Original]} Authentication is the process of:
  \begin{enumerate}[label=\Alph*.]
    \item granting permission to a verified user.
    \item verifying that a user or system is who or what it claims to be.
    \item encrypting a message before sending it.
    \item backing up a database.
  \end{enumerate}
\item \textit{[Original]} Which statement correctly distinguishes
  \textbf{authentication} from \textbf{authorization}?
  \begin{enumerate}[label=\Alph*.]
    \item Authentication grants permission; authorization proves identity.
    \item Authentication proves identity; authorization grants permission.
    \item Both authentication and authorization prove identity.
    \item Both authentication and authorization grant permission.
  \end{enumerate}
\item \textit{[Original]} Expand OTP.
  \begin{enumerate}[label=\Alph*.]
    \item One-Time Password
    \item Online Transfer Password
    \item Open Transaction Protocol
    \item One-Time Program
  \end{enumerate}
\item \textit{[Original]} Two-factor authentication means using:
  \begin{enumerate}[label=\Alph*.]
    \item two passwords of different lengths.
    \item two different authentication factors.
    \item two separate computers to log in.
    \item two firewalls on the same network.
  \end{enumerate}
\item \textit{[Original]} A fingerprint used for login is an example of which
  authentication factor?
  \begin{enumerate}[label=\Alph*.]
    \item Something you know
    \item Something you have
    \item Something you are
    \item Something you forget
  \end{enumerate}
\item \textit{[Original]} Which of the following is an example of a
  \textbf{digital public service}?
  \begin{enumerate}[label=\Alph*.]
    \item Applying for a certificate through an online government portal
    \item Buying goods from a private online shop
    \item Sending a personal email to a friend
    \item Using a word processor to type a letter
  \end{enumerate}
\item \textit{[Original]} Who is responsible for managing user accounts and
  access rights on a computer system?
  \begin{enumerate}[label=\Alph*.]
    \item The end user
    \item The system administrator
    \item The data-entry operator
    \item The programmer
  \end{enumerate}
\item \textit{[Original]} Which pairing is correctly matched?
  \begin{enumerate}[label=\Alph*.]
    \item User --- writes and maintains the software
    \item Operator --- runs the computer system and routine jobs
    \item Programmer --- manages the network
    \item Technician --- uses the finished system to do office work
  \end{enumerate}
\end{enumerate}
\end{document}
